\documentclass[10pt,english,fleqn,openany,twoside,letterpaper]{book}

%Muestra los márgenes del documento para evitar Warnings
%Para activar la siguiente línea quite el simbolo % 
%\usepackage[showframe]{geometry}

%Formato de fuentes bibliográficas
%Use el estilo bibliográfico que sea pertinente según el área de estudio APA, IEEE, etc

%Usando el paquete BibLaTeX
%Cita normal con \cite[página]{} y cita con paréntesis \parencite[página]{}

% Configuración de BibLaTeX
%\usepackage[backend=biber,style=authoryear,maxcitenames=2,maxbibnames=99,giveninits=true,uniquename=false]{biblatex}
%\addbibresource{biblio.bib}

% Cambiar el idioma de las referencias bibliográficas a español
%\DefineBibliographyStrings{spanish}{%
%  andothers = {et\addabbrvspace al\adddot},
%  andmore = {et\addabbrvspace al\adddot},
%}

% Personalizar el formato de las citas y la bibliografía
%\DeclareNameAlias{sortname}{family-given}
%\DeclareDelimFormat{multinamedelim}{\addcomma\space}
%\DeclareDelimFormat{finalnamedelim}{\addcomma\space\&\space}
%\DeclareFieldFormat{titlecase}{\MakeSentenceCase*{#1}}
%\DeclareFieldFormat[article,inbook,incollection,inproceedings,patent,thesis,unpublished]{title}{\titlecase{#1}}
%\DeclareFieldFormat{journaltitlecase}{\titlecase{#1}}
%\DeclareFieldFormat{pages}{#1}
%\DeclareFieldFormat{volume}{\mkbibbold{#1}}
%\renewbibmacro{in:}{}
%\AtEveryBibitem{\clearfield{month}}

%Usando el paquete Natbib
%Cita normal \cite[página]{} y cita con paréntesis \citep[página]{}
\usepackage[numbers]{natbib}
\bibpunct{[}{]}{;}{\&}{.}{}
\bibliographystyle{unsrt}

%Idioma del documento
%Use main para el idioma principal del documento
\usepackage[main=english,spanish]{babel}

% Carácteres especiales
\usepackage{fontenc}
\usepackage{lmodern}

% Evita ligadura li & fl
\usepackage{microtype}
\DisableLigatures{encoding = *, family = *}

% Otros paquetes de tablas y colores avanzados
\usepackage{amsmath,graphicx,rotating,float,multirow}
\usepackage{longtable}
\setlength{\LTcapwidth}{6in}
\usepackage[utf8]{inputenc}
\usepackage{epsfig,epic,eepic,threeparttable,amscd,here,lscape,tabularx,subfigure}
\usepackage{tabu,array}
\usepackage[rgb]{xcolor}

% Permite ver y configurar los parámetros de la página
\usepackage{layout}
%Hyperref permite ver las secciones del texto
\usepackage[hidelinks]{hyperref}

%Permite incluir código de cualquier lenguaje dentro del texto del documento
\usepackage{minted}
\usepackage{fancyvrb}
\newenvironment{myverbatim}{\Verbatim}{\endVerbatim}

%Genera los comandos de la página de autoría
\newcommand{\studentname}{}
\newcommand{\submissiondate}{}
\newcommand{\academictitle}{}
\newcommand{\resgroupone}{}
\newcommand{\resgrouptwo}{}
\newcommand{\researchtopic}{}
\newcommand{\thesisname}{}
\newcommand{\thesisnameeng}{}
\newcommand{\thesisnamelang}{} %Usar solo si se requiere
\newcommand{\director}{}
\newcommand{\directortitle}{}
\newcommand{\codirector}{} %Usar solo si se requiere
\newcommand{\codirectortitle}{} %Usar solo si se requiere
\newcommand{\issuedate}{}

\newcommand{\keywords}{}

\newcommand{\sede}{}
\newcommand{\department}{}
\newcommand{\departmenttwo}{} %Usar solo si se requiere
\newcommand{\faculty}{}
\newcommand{\university}{Universidad Nacional de Colombia}

%Información de la tesis
%Diligenciar aquí los datos para su carga automática donde se requiera en el documento
%En el caso de tesis o trabajos finales, verificar que el título coincida con el aprobado por la Facultad
\renewcommand{\studentname}{Wilson Andrés Pinillo Ipia}
\renewcommand{\thesisname}{Optimización Dinámica de Precios mediante Aprendizaje Reforzado
Profundo}
\renewcommand{\thesisnameeng}{Dynamic Price Optimization using Deep Reinforcement Learning}
%\renewcommand{\thesisnamelang}{Nombre del trabajo o tesis en un tercer idioma} %Usar solo si se requiere
\renewcommand{\issuedate}{2025}
\renewcommand{\submissiondate}{Fecha entrega}
\renewcommand{\director}{Ph.D. John Willian Branch Bedoya}
\renewcommand{\directortitle}{Profesor Titular}
%\renewcommand{\codirector}{Prof. Dr. Co director}
%\renewcommand{\codirectortitle}{Indicar si es Profesor Titular/Asociado}
\renewcommand{\academictitle}{Magíster en Ingeniería - Analítica}
\renewcommand{\resgroupone}{Grupo A (Sigla Grupo Investigación 01) }
\renewcommand{\resgrouptwo}{Grupo B (Sigla Grupo Investigación 02) }
\renewcommand{\researchtopic}{Línea}
\renewcommand{\sede}{Sede Medellín} 
\renewcommand{\department}{Departamento de Ciencias de la Computación y de la Decisión}
\renewcommand{\departmenttwo}{Departamento 2} %Usar solo si es necesario
\renewcommand{\faculty}{Facultad de Minas}

%Palabras clave del documento - Tener presente los Theasurus https://www.thesaurus.com/
%Disponible en 3 idiomas aunque se puede extender a francés o otro idioma
%\renewcommand{\schlusselworter}{}
%\renewcommand{\palavraschave}{}

% Estilo de los encabezados y pies de página
\usepackage{fancyhdr}
\fancyhf{}%
\pagestyle{fancyplain}
\textheight22.5cm \topmargin0cm \textwidth16.5cm \headheight22pt
\oddsidemargin0.5cm \evensidemargin-0.5cm%
\fancypagestyle{plain}{
\fancyhead[RO,LE]{}
\fancyhead[RE,LO]{\small \textbf{\thesisnameeng}}
\fancyfoot[CO,CE]{\thepage}
}
\pagestyle{fancy}
\fancyhf{}%
\renewcommand{\chaptermark}[1]{\markboth{\thechapter.\; #1}{}}
\renewcommand{\sectionmark}[1]{\markright{\thesection.\; #1}{}}
\fancyhead[LO,RE]{\leftmark}
\fancyhead[RO,LE]{\rightmark}
\fancyfoot[CO,CE]{\thepage}
\thispagestyle{fancy}%

\usepackage{titlesec}
% Permite personalizar los títulos de sección y de capítulos
% hang lo deja en el mismo renglón, display lo despliega
% Elimina el "Capitulo" y deja solo el número
\titleformat{\chapter}[hang]
  {\sffamily\Huge\bfseries}{\thechapter}{0.5cm}{\sffamily\Huge}
\titleformat{\section}[hang]{\sffamily\LARGE}{\thesection}{0.5cm}{}
\titleformat{\subsection}[hang]{\sffamily\Large}{\thesubsection}{0.5cm}{}
\titleformat{\subsubsection}[hang]{\sffamily\large}{\thesubsubsection}{0.5cm}{}
\titleformat{\paragraph}[runin]{\sffamily\normalsize}{}{}{\emph}

%Coloca anexo o apéndice en la Tabla de contenido
\usepackage[toc,page]{appendix}

% Configuración de las páginas en twoside-mode
% Permite ver y configurar los parámetros de la página
\setlength{\voffset}{-0.25in}
\setlength{\headwidth}{467pt}
\setlength{\headheight}{22pt}
\setlength{\oddsidemargin}{0pt}
\setlength{\evensidemargin}{0pt}
\setlength{\marginparwidth}{0pt}
\setlength{\marginparsep}{0pt}
\setlength{\parskip}{2em}
\setlength{\footskip}{20pt}
\setlength{\textheight}{650pt}
\setlength{\textwidth}{467pt}
\setlength{\headsep}{5pt}
\setlength{\parindent}{0pt}
\setlength{\baselineskip}{10pt plus 5pt minus 5pt}
\renewcommand{\theequation}{\thechapter-\arabic{equation}}
\renewcommand{\thefigure}{\textbf{\thechapter-\arabic{figure}}}
\renewcommand{\thetable}{\textbf{\thechapter-\arabic{table}}}

%Ajusta el espacio entre la etiqueta de figuras y tablas y su título en la lista de figuras y en la de tablas 
\usepackage{titletoc} 
\titlecontents{figure}[0em]{}{\thecontentslabel\hspace{1em}}{}{\titlerule*[1pc]{.}\contentspage}
\titlecontents{table}[0em]{}{\thecontentslabel\hspace{1em}}{}{\titlerule*[1pc]{.}\contentspage}

%Define la distancia de la primera linea de un parrafo a la margen
\parindent0cm 

%Espacio entre lineas
\renewcommand{\baselinestretch}{1}

%Permite personalizar el ajuste vertical mediante cajas
\usepackage{adjustbox}

%Para rotar texto, objetos, tablas y páginas.
\usepackage{rotating}

%Permite incluir mecanismos y reacciones químicas
\usepackage{tikz}
\usepackage{chemformula}
\usepackage{chemfig}

\usetikzlibrary{calc,arrows.meta}% per right to e left to
\tikzset{
myedge/.style={->, -{Latex[#1]}}
}

%Fuente de la presentación Ancizar Sans UNAL
%Para usar este compilado en Overleaf se debe usar el compilador XeLaTeX o LuaLaTeX!!
%Menu -> Compiler -> XeLaTeX o LuaLaTeX
%La siguiente línea debe comentarse si desea compilar con pdfLaTeX
%\RequireXeTeX

% Definición de la fuente Ancizar Sans
\newif\ifxetexorluatex

\ifxetexorluatex
  \usepackage{fontspec}
  \usefonttheme{serif}
  \setmainfont{AncizarSans}[Path=./AncizarSans/,Scale=1,Extension=.otf,UprightFont=*-Regular,BoldFont=*-Bold,ItalicFont=*-Italic,BoldItalicFont=*-BoldItalic]
\else
  % Si se compila con pdfLaTeX, cargar la fuente apropiada aquí
  \usepackage[T1]{fontenc}
\fi
% Metadatos del documento
\AtBeginDocument{%
	\hypersetup{
		pdfborder={0 0 0},
		pdfauthor={\studentname},
		pdfsubject={\thesisname}, 
		pdfcreator={\studentname},
		pdfproducer={\studentname},
	}
}

%Carga el simbolo de grado y el de Angstrom
\newcommand{\angstrom}{\textup{\AA}}
\newcommand{\grad}{$^{\circ}$}

%Inicio del documento, no olvide la etiqueta de cierre al final \end{document}
\begin{document}

%Nombres y formatos de títulos, tablas y figuras
%Use \sffamily para dejar con letra Sans Serif, sin etiqueta queda LaTeX clásico
\renewcommand{\listfigurename}{\sffamily List of Figures}
\renewcommand{\listtablename}{\sffamily List of Tables}
\renewcommand{\contentsname}{\sffamily Contents}
\renewcommand{\chaptername}{\sffamily Chapter}
\renewcommand{\tablename}{\scriptsize \centering \textbf{Table}}
\renewcommand{\figurename}{\scriptsize \centering \textbf{Figure}}
\renewcommand{\appendixname}{\sffamily Appendix}

%Cambia el nombre de la sección de referencias
\renewcommand{\bibname}{\sffamily References}

%Páginas de Presentación del documento - No modificar esto se hace automáticamente
{\newpage
\thispagestyle{empty}
\begin{center}
\begin{figure}
\centering
\epsfig{file=00Figuras/00f00EscudoUN2016,scale=1}%
\end{figure}
\textbf{\Huge \thesisnameeng} \\ 
\vspace{2.5cm}
\textbf{\Large \studentname} \\
\vspace{5.0cm}
\faculty \\ \department \\
\sede, Colombia \\
\issuedate
\newpage 
\thispagestyle{empty}
\vspace{2.0cm}
%\small Nota: No utilizar el tipo de letra Ancizar en el documento puesto que este tipo de fuente restringe la copia de los archivos en el Repositorio Institucional. [Recuerde borrar esta nota]\\
\textbf{\Huge \thesisnameeng} \\
\vspace{2.0cm}
\textbf{\Large \studentname} \\
\vspace{2.0cm}
\small Final project presented as a partial requirement to apply for the title of: \\
{\bfseries \academictitle}\\
\vspace{2.0cm}
\textbf{Director:} \\
\director \\
\directortitle \; - \departmenttwo \\
\faculty \\
\university \\ 
\vspace{0.5cm}
%\textbf{Co-director:} \\
%\codirector \\
%\codirectortitle \; - \department \\
%\faculty \\
%\university 
%\vspace{1.5cm} \\
\textbf{Research line:} \\ 
\researchtopic\\
\textbf{Research group:} \\
\resgroupone \\
\resgrouptwo \\
\vspace{1.5cm} 
\university \\
\faculty \\
\department \\
\issuedate
\end{center}

% Dedicatorias
\newpage
\thispagestyle{empty}
\begin{flushright}
\begin{minipage}{12.5cm}
\noindent
\\[10em]
%Modificar la cita que se quiere agregar
{\Large Cita 01.}
\\[3em]
Autor
\\ \textit{Fuente}
\\[10em]
%Para anular la adición de una segunda cita anule las siguientes lineas desde acá mediante comentario (%)
{\Large \textit{Wenn du es nicht einfach erkl\"{a}ren kannst, hast du es nicht genug verstanden} - Si no eres capaz de explicar algo claramente, es que aún no lo has entendido lo suficiente.}
\\[3em]
Albert Einstein
%Hasta acá!
\end{minipage}
\end{flushright} 


% Declaracíon de originalidad del texto y del contenido
% No modificar, se hace automáticamente con los comandos ya definidos
\newpage
\chapter*{\sffamily Declaration}
\par I hereby declare that I have prepared this thesis independently, and that I have not used any sources or aids other than those specified. All passages taken literally or in spirit from published or unpublished texts have been identified as such. No part of this work has been used in any other thesis.
\\[1em]
\sede., \submissiondate
\\[6em]
\rule{6cm}{0.5pt}\\
\studentname
}

%Páginas preámbulo, listado de figuras, tablas y tabla de contenido
{\pagestyle{plain} \pagenumbering{roman}
\setlength{\parskip}{1mm}
\chapter*{\sffamily Acknowledgments}
\addcontentsline{toc}{chapter}{Acknowledgments}%
% Comentar las dos lineas de abajo con % en caso que no se requieran abreviaturas y resumen en el trabajo
\chapter*{\sffamily List of Symbols and Abbreviations}
\addcontentsline{toc}{chapter}{List of Symbols and Abbreviations}
%\newpage
%\thispagestyle{empty}
%\chapter*{\sffamily Resumen}
%\addcontentsline{toc}{chapter}{Resumen}%
%\begin{center}
%\textbf{\large \thesisname}
%\end{center}
%\par Texto del resumen.
%\\[2cm]
%\textbf{Palabras clave:} \palabrasclave

\newpage
\thispagestyle{empty}
\chapter*{\sffamily Abstract}
\addcontentsline{toc}{chapter}{Abstract}%
\begin{center}
\textbf{\large \thesisnameeng}
\end{center}
\par Abstract text.
\\[2cm]
\textbf{Keywords:} \keywords

%Usar si es requerido, de no necesitarse favor comentar con % las lineas
%\newpage 
%\thispagestyle{empty}
%\chapter*{\sffamily Zusammenfassung}
%\addcontentsline{toc}{chapter}{Zusammenfassung}%
%\begin{center}
%\textbf{\large \thesisnamelang}
%\end{center}
%\par Zusammenfassung Texte.
%\\[2cm]
%\textbf{Schlüsselwörter:} \schlusselworter
% Dejar esta parte así para que genere correctamente la página de la tabla de contenido
\addcontentsline{toc}{chapter}{List of Figures}
\listoffigures
\clearpage
\addcontentsline{toc}{chapter}{List of Tables}
\listoftables
\clearpage
\addcontentsline{toc}{chapter}{Contents}
\tableofcontents
\clearpage
}

{\pagenumbering{arabic}
\setlength{\parskip}{\baselineskip}
%Incluir secciones del documento de aqui en adelante
%Use \include para incluir desde una página nueva e \input para incluir sin salto de página
\chapter{Problem Statement}
In the retail industry, the rapid expansion of e-commerce platforms and real-time price comparison tools has intensified competition, making pricing decisions increasingly critical to maintaining profitability and competitiveness. Many retailers still rely on static or rule-based pricing strategies that lack the flexibility to adapt to complex, rapidly changing market dynamics. These limitations often result in suboptimal revenue and margin performance, particularly in highly competitive and price-sensitive markets \cite{bauer2018optimal, zhu2024drl}. Furthermore, despite the growing availability of extensive transactional and customer-level data, most retailers fail to effectively leverage these resources to develop adaptive, data-driven pricing strategies. Recent advances in reinforcement learning have demonstrated strong potential for sequential decision-making under uncertainty, providing a framework well suited to pricing problems that require continuous adaptation to fluctuating demand and evolving market conditions \cite{cai2022survey, qiao2024distributed}. Nevertheless, the application of Deep Reinforcement Learning (DRL) to dynamic price optimization remains limited, despite its ability to exploit large-scale retail data for adaptive policy learning \cite{liu2019dynamic}. This gap underscores the need to design, implement, and evaluate a DRL-based framework capable of learning optimal pricing policies through interaction with a simulated retail environment, with the objective of maximizing profitability and operational efficiency \cite{zhu2024drl}.

\chapter{Objectives}
\include{01Section01}
\chapter{Theoretical Framework}
\chapter{Materials and Methods}

% TEMPORARY NOTE: Explanation of DQNPolicy
% This section will be expanded later.

\subsection*{Technical Note: DQNPolicy Implementation in Stable-Baselines3}

The \texttt{DQNPolicy} class is central to the agent's learning process. It manages the neural networks that approximate the Q-function, which estimates the value of taking an action in a given state. Its operation can be summarized in a few key points:

\begin{itemize}
    \item \textbf{Two-Network Architecture:} The policy initializes two identical neural networks:
    \begin{itemize}
        \item \textbf{Online Network (\texttt{q\_net}):} This network is used for action selection during interaction with the environment. Its weights are continuously updated during training.
        \item \textbf{Target Network (\texttt{q\_net\_target}):} This is a periodically updated, frozen copy of the online network. It is used to provide a stable target for the learning updates, which is a crucial technique to prevent oscillations and stabilize the training process.
    \end{itemize}

    \item \textbf{Action Selection:} To choose an action, the current state is fed into the online network, which outputs a Q-value for every possible action. The policy then selects the action with the highest Q-value (a greedy policy).

    \item \textbf{Learning Process:} The learning occurs by sampling a batch of experiences from a replay buffer. The online network's weights are updated to minimize the loss between its predicted Q-values and the target Q-values calculated using the target network. Specifically, the target is computed as:
    \[ y = r + \gamma \max_{a'} Q_{\text{target}}(s', a') \]
    The loss is then typically the Mean Squared Error between this target $y$ and the prediction from the online network, $Q_{\text{online}}(s, a)$.
\end{itemize}

% TEMPORARY NOTE: Explanation of PPO
% This section will be expanded later.

\subsection*{Technical Note: PPO Implementation in Stable-Baselines3}

Proximal Policy Optimization (PPO) is an \textbf{Actor-Critic} algorithm that has become a go-to choice for many reinforcement learning problems due to its reliable performance and ease of implementation. Unlike value-based methods like DQN, PPO learns a policy directly.

Its operation can be summarized as follows:

\begin{itemize}
    \item \textbf{Actor-Critic Architecture:} PPO uses two distinct neural networks:
    \begin{itemize}
        \item \textbf{The Actor:} This network represents the policy. It takes the current state as input and outputs a probability distribution over the possible actions. It is responsible for deciding which action to take.
        \item \textbf{The Critic:} This network estimates the state-value function, $V(s)$. Its job is not to choose actions, but to evaluate how good a state is, which in turn helps the Actor learn more efficiently by providing a low-variance estimate of performance.
    \end{itemize}

    \item \textbf{The PPO Objective Function:} The core innovation of PPO is its objective function, which prevents the policy from changing too drastically at each update step. This is achieved with a "clipped surrogate objective".
    \begin{enumerate}
        \item It calculates the probability ratio between the new policy and the old policy: $r_t = \pi_{\text{new}}(a|s) / \pi_{\text{old}}(a|s)$.
        \item It uses the Critic to compute an \textbf{Advantage} value ($\hat{A}_t$), which represents how much better an action was than the average action in that state.
        \item The policy update is then "clipped" to keep the new policy close to the old one, avoiding destructive updates. The simplified objective is:
        \[ L_{\text{CLIP}} = \mathbb{E} \left[ \min \left( r_t \hat{A}_t, \text{clip}(r_t, 1-\epsilon, 1+\epsilon) \hat{A}_t \right) \right] \]
    \end{enumerate}

    \item \textbf{Learning Process:} PPO's training loop is different from DQN's.
    \begin{itemize}
        \item The Actor interacts with the environment for a set number of steps to collect a batch of experiences.
        \item Using this batch, the Critic helps compute the advantage for each step.
        \item The algorithm then performs multiple epochs of gradient ascent on the collected data, updating both the Actor (using the clipped objective) and the Critic (by minimizing the error in its value predictions). This ability to reuse experience for multiple updates makes it more data-efficient than simpler policy gradient methods.
    \end{itemize}
\end{itemize}
\chapter{Results}
\chapter{Discussion of Results}
\chapter{Conclusions}
\chapter{Recommendations}

%Inicio del apéndice o anexos
\begin{appendix}
\chapter{Appendix 1}%
\end{appendix}

%Permite visualizar la bibliografía en la tabla de contenido
%Cambie el nombre a Bibliografía o Literatura Citada en la siguiente línea de ser preciso
\addcontentsline{toc}{chapter}{References} 

\let\OLDthebibliography=\thebibliography
\def\thebibliography#1{\OLDthebibliography{#1}}
{\scriptsize
\pagestyle{plain}
% Nombre del documento donde se almacenan las referencias
\bibliography{References}
\nocite{*}
% Inserta un página adicional al final en blanco 
%\cleardoublepage
% Para NO insertar una página adicional al final usar \clearpage
\clearpage
}}

\end{document}